This chapter provides brief guidance on interpreting the references cited in this thesis. All references follow the IEEE citation format as required by the SOICT graduation thesis template. The full reference list is automatically generated from the \texttt{reference.bib} BibTeX database.

References are cited inline using the \texttt{\textbackslash cite\{key\}} command, which produces bracketed numbers (e.g., \cite{nakamoto2008bitcoin}). Where multiple sources support the same claim, they are cited together (e.g., \cite{buterin2014ethereum, yakovenko2018solana}).

The primary references consulted in this thesis fall into five categories:

\begin{enumerate}
    \item \textbf{Blockchain foundational works:} The Bitcoin whitepaper \cite{nakamoto2008bitcoin}, the Ethereum whitepaper \cite{buterin2014ethereum}, and the Solana architecture paper \cite{yakovenko2018solana} provide the theoretical foundation for the blockchain layer.
    \item \textbf{Smart contract and DApp development:} The Anchor framework documentation \cite{anchor2021} and Antonopoulos and Wood's \textit{Mastering Ethereum} \cite{antonopoulos2018mastering} informed the smart contract design patterns.
    \item \textbf{Decentralized identity and reputation:} The Decentralized Society paper by Buterin, Weyl, and Ohlhaver \cite{weyl2022decentralized} provides the theoretical basis for Soulbound Tokens.
    \item \textbf{Industry reports:} Statista's freelance market report \cite{statista2024freelance} and the annual reports of Upwork \cite{upwork2023} and Fiverr \cite{fiverr2023} provide market context.
    \item \textbf{Technology documentation:} Official documentation for React \cite{react2023}, Next.js \cite{nextjs2023}, Material UI \cite{materialui2023}, IPFS \cite{ipfs2015}, the Meta LLaMA 3 technical report \cite{touvron2023llama}, and the Groq API documentation \cite{groq2024api} informed the frontend and AI integration design.
\end{enumerate}
